\section{Algoritmo de Sobrecarga Progresiva Automática}

\subsection{Formulación Matemática de Estimación de 1RM}

Para cuantificar la fuerza relativa en una serie de $r$ repeticiones ejecutadas con carga $w$ (expresada en $kg$) a una percepción de esfuerzo $RPE \in [6.0, 10.0]$, calculamos en primer lugar las repeticiones en recámara equivalentes (\textit{Repetitions in Reserve}, $RIR$):
\begin{equation}
    RIR = 10.0 - RPE
\end{equation}

Las repeticiones estimadas a fallo absoluto $r_{\text{max}}$ se determinan como:
\begin{equation}
    r_{\text{max}} = r + RIR
\end{equation}

Utilizamos una combinación ponderada de los modelos no lineales de \textbf{Epley} y \textbf{Brzycki}:

\begin{enumerate}
    \item \textbf{Fórmula de Epley:}
    \begin{equation}
        1RM_{\text{Epley}} = w \cdot \left(1 + \frac{r_{\text{max}}}{30}\right)
    \end{equation}

    \item \textbf{Fórmula de Brzycki:}
    \begin{equation}
        1RM_{\text{Brzycki}} = w \cdot \frac{36}{37 - r_{\text{max}}}
    \end{equation}
\end{enumerate}

El 1RM estimado final de la serie $1RM_{\text{est}}$ se obtiene promediando ambos modelos:
\begin{equation}
    1RM_{\text{est}} = \frac{1RM_{\text{Epley}} + 1RM_{\text{Brzycki}}}{2}
\end{equation}

\subsection{Volumen Efectivo e Incremento Progresivo Automático}

El \textbf{Volumen Efectivo Total} ($V_{\text{efectivo}}$) de una sesión para un ejercicio dado, considerando únicamente series efectivas ($is\_warmup = 0$), se expresa como:
\begin{equation}
    V_{\text{efectivo}} = \sum_{i=1}^{N} w_i \cdot r_i \cdot \alpha(RPE_i)
\end{equation}

donde $\alpha(RPE_i)$ es el coeficiente de ponderación de fatiga:
\begin{equation}
    \alpha(RPE) = \begin{cases}
        1.0 & \text{si } RPE \ge 8.0 \\
        0.85 & \text{si } 7.0 \le RPE < 8.0 \\
        0.70 & \text{si } RPE < 7.0
    \end{cases}
\end{equation}

\subsection{Regla de Incremento para la Siguiente Sesión}

La recomendación de carga $w_{\text{next}}$ para la subsiguiente sesión se ajusta según la media de $RIR_{\text{prom}}$ de la sesión actual y el logro del límite superior del rango de repeticiones objetivo $[R_{\text{min}}, R_{\text{max}}]$:

\begin{equation}
    w_{\text{next}} = \begin{cases}
        w_{\text{actual}} + \Delta w & \text{si } r_i \ge R_{\text{max}} \;\forall i \quad \text{y} \quad RIR_{\text{prom}} \ge 1.5 \\
        w_{\text{actual}} & \text{si } R_{\text{min}} \le r_i < R_{\text{max}} \quad \text{o} \quad 0.5 \le RIR_{\text{prom}} < 1.5 \\
        w_{\text{actual}} \cdot (1 - \beta) & \text{si } RIR_{\text{prom}} < 0.5 \quad \text{(deload o fatiga excesiva)}
    \end{cases}
\end{equation}

donde $\Delta w$ es el incremento estándar ($2.5\,kg$ para ejercicios compuestos, $1.25\,kg$ para ejercicios de aislamiento) y $\beta \in [0.05, 0.10]$ es el factor de descarga por fatiga acumulada.
